\documentclass[11pt,a4paper]{scrartcl}

\usepackage[lithuanian]{babel}
\usepackage{../../1.\space Vienanariai\space ir\space daugianariai/manostilius_lt}
\usepackage{tasks}
\renewcommand{\theproblem}{\arabic{problem}}

\title{Kvadratinio trinario skaidymas dauginamaisiais}
\author{Andrius Berniukevičius}

\begin{document}

\maketitle
\section*{Kaip dirbti su uždaviniais?}

Pirmiausia kiekvieną uždavinį pabandykite išspręsti savarankiškai. Neskubėkite ieškoti užuominų — skirkite laiko apgalvoti, ką jau žinote ir kokį sprendimo būdą galėtumėte pritaikyti. Jei įstrigote, atsiverskite atitinkamą „Hintai 1“ užuominą. Pagalvojate. Jei jos nepakanka, tik tada skaitykite „Hintai 2“, o prireikus — „Hintai 3“. Užuominas naudokite palaipsniui, nežiūrėkite jų iš anksto.

Paties pagimdytas sprendimas yra šimtą kartų vertingesnis už įvaikintą sprendimą.

\section{Uždaviniai}

\begin{problem}
Iškelkite bendrą dauginamąjį prieš skliaustus:
\begin{tasks}[label={(\arabic*)}, label-offset=0.9em](2)
    \task $3x + 12$
    \task $5ab + 10ac$
    \task $x^2 + 7x$
    \task $4x^2 - 12x$
    \task $-2x + 10$
    \task $-x_2x + x_1x_2$
\end{tasks}
\end{problem}

\begin{problem}
Iškelkite bendrą dvinario dauginamąjį prieš skliaustus:
\begin{tasks}[label={(\arabic*)}, label-offset=0.9em](2)
    \task $3x(x + 2) + 5(x + 2)$
    \task $y(2y - 1) - 4(2y - 1)$
    \task $a(a - 3) + (a - 3)$
    \task $2m(m + 4) - 3n(m + 4)$
    \task $x(x - 5) - (x - 5)$
    \task $-3x(x - 4) - 2(4-x)$
    \task $ab(b-2c) - c(2c-b)$
    \task $a(b-a)+b(a-b)-c(a-b)$
    \task $x(y-z) + y(z-y)-xyz(y-z)$
    \task $ab(abc-1)+bc(abc-1)-ca(abc-1)$
    \task $x(x-x_1)-x_2(x-x_1)$  
    \task $x(x-x_1)-x_2(x_1-x)$  
\end{tasks}
\end{problem}

\begin{problem}
Skaidykite dauginamaisiais grupuodami narius poromis t.y. iš pradžių iškeldami bendrą dauginamąjį iš pirmos poros narių, o tada iš antros poros narių:
\begin{tasks}[label={(\arabic*)}, label-offset=0.9em](2)
    \task $x^2 + 2x + 3x + 6$
    \task $2a^2 - 6a + 5a - 15$
    \task $6m^2 + 2m + 9m + 3$
    \task $y^2 - 4y - y + 4$
    \task $x^2 - 7x - 2x + 14$
    \task $ab^{2}+abc+bc+ c^{2} $ 
    \task $x^3 + 3x^2 + 2x + 6$
    \task $x ^{3} -x ^{2} -y ^{2} x  +y ^{2} $ 
\end{tasks}
\end{problem}

\begin{problem}
Duotas kvadratinis trinaris $ax^2+bx+c$, kur $a\ne 0$. Tegul $x_1$ ir $x_2$ yra lygties $ax^2+bx+c=0$ šaknys. Įrodykite, kad trinarį galima užrašyti pavidalu:
\[
ax^2+bx+c=a(x-x_1)(x-x_2).
\]
\end{problem}

\begin{problem}
Naudodamiesi kvadratinio trinario skaidymo dauginamaisiais formule išskaidykite šiuos kvadratinius trinarius:
\begin{tasks}[label={(\arabic*)}, label-offset=0.9em](2)
    \task $x^2 - 6x + 8$
    \task $6x^2 + 5x - 1$
    \task $x^2 - x - 12$
    \task $10x^2 + 3x -1$
    \task $x^2 - 9x + 14$
    \task $-x^2 -3x +18$
    \task $2x^2 - 14x + 24$
    \task $-3x^2 - 6x + 9$
    \task $2x^2 + 5x - 3$
    \task $-5x^2 + 3x + 2$
    \task $3x^2 - 7x - 6$
    \task $-2x^2 + x + 3$
\end{tasks}
\end{problem}

\begin{problem}
Išskaidykite kvadratinį trinarį ir suprastinkite algebrines trupmenas:
\begin{tasks}[label={(\arabic*)}, label-offset=0.9em](2)
    \task $\frac{x^2 - 3x - 10}{x - 5}$
    \task $\frac{x^2 - 9}{x^2 - x - 6}$
    \task $\frac{2x^2 + x - 1}{x + 1}$
    \task $\frac{3x^2 - 12x}{x^2 - 2x - 8}$
    \task $\frac{x^2 - 4x - 12}{x - 6}$
    \task $\frac{x^2 - 4}{x^2 + x - 6}$
\end{tasks}
\end{problem}

\newpage
\section{Atsakymai}

\begin{enumerate}[label=\textbf{\arabic*.}, leftmargin=*]
    \item
    \begin{enumerate}[label={(\arabic*)}]
        \item $3(x+4)$
        \item $5a(b+2c)$
        \item $x(x+7)$
        \item $4x(x-3)$
        \item $-2(x-5)$
        \item $x_2(x_1-x)$
    \end{enumerate}
    \item
    \begin{enumerate}[label={(\arabic*)}]
        \item $(x+2)(3x+5)$
        \item $(2y-1)(y-4)$
        \item $(a-3)(a+1)$
        \item $(m+4)(2m-3n)$
        \item $(x-5)(x-1)$
        \item $(x-4)(2-3x)$
        \item $(b-2c)(ab+c)$
        \item $(a-b)(b-a-c)$
        \item $(y-z)(x-y-xyz)$
        \item $(abc-1)(ab+bc-ca)$
        \item $(x-x_1)(x-x_2)$
        \item $(x-x_1)(x+x_2)$
    \end{enumerate}
    \item
    \begin{enumerate}[label={(\arabic*)}]
        \item $(x+2)(x+3)$
        \item $(a-3)(2a+5)$
        \item $(3m+1)(2m+3)$
        \item $(y-4)(y-1)$
        \item $(x-7)(x-2)$
        \item $(b+c)(ab+c)$
        \item $(x+3)(x^2+2)$
        \item $(x-1)(x-y)(x+y)$
    \end{enumerate}
    \setcounter{enumi}{4}
    \item
    \begin{enumerate}[label={(\arabic*)}]
        \item $(x-2)(x-4)$
        \item $(6x-1)(x+1)$
        \item $(x+3)(x-4)$
        \item $(5x-1)(2x+1)$
        \item $(x-7)(x-2)$
        \item $(3-x)(x+6)$
        \item $2(x-3)(x-4)$
        \item $-3(x-1)(x+3)$
        \item $(2x-1)(x+3)$
        \item $(1-x)(5x+2)$
        \item $(3x+2)(x-3)$
        \item $(3-2x)(x+1)$
    \end{enumerate}
    \item
    \begin{enumerate}[label={(\arabic*)}]
        \item $x+2$, kai $x\ne5$
        \item $\frac{x+3}{x+2}$, kai $x\ne3,-2$
        \item $2x-1$, kai $x\ne-1$
        \item $\frac{3x}{x+2}$, kai $x\ne4,-2$
        \item $x+2$, kai $x\ne6$
        \item $\frac{x+2}{x+3}$, kai $x\ne-3,2$
    \end{enumerate}
\end{enumerate}

\newpage
\section{Hintai 1}

\begin{enumerate}[label=\textbf{\arabic*.}, leftmargin=*]
    \item Iškelkite bendrą skaitinį arba raidinį daugiklį.
    \begin{enumerate}[label={(\arabic*)}]
        \item Raskite bendrąjį daugiklį reiškinyje $3x+12$.
        \item Raskite bendrąjį daugiklį reiškinyje $5ab+10ac$.
        \item Raskite bendrąjį daugiklį reiškinyje $x^2+7x$.
        \item Raskite bendrąjį daugiklį reiškinyje $4x^2-12x$.
        \item Raskite bendrąjį daugiklį reiškinyje $-2x+10$; atkreipkite dėmesį į ženklą.
        \item Raskite bendrąjį daugiklį reiškinyje $-x_2x+x_1x_2$.
    \end{enumerate}
    \item Ieškokite abiejuose nariuose pasikartojančio dvinario.
    \begin{enumerate}[label={(\arabic*)}]
        \item Raskite bendrąjį dvinario daugiklį reiškinyje $3x(x+2)+5(x+2)$.
        \item Raskite bendrąjį dvinario daugiklį reiškinyje $y(2y-1)-4(2y-1)$.
        \item Raskite bendrąjį dvinario daugiklį reiškinyje $a(a-3)+(a-3)$.
        \item Raskite bendrąjį dvinario daugiklį reiškinyje $2m(m+4)-3n(m+4)$.
        \item Raskite bendrąjį dvinario daugiklį reiškinyje $x(x-5)-(x-5)$.
        \item Raskite bendrąjį dvinario daugiklį, atsižvelgdami į $4-x=-(x-4)$.
        \item Raskite bendrąjį dvinario daugiklį, atsižvelgdami į $2c-b=-(b-2c)$.
        \item Raskite bendrąjį dvinario daugiklį, atsižvelgdami į $b-a=-(a-b)$.
        \item Raskite bendrąjį dvinario daugiklį, atsižvelgdami į $z-y=-(y-z)$.
        \item Raskite bendrąjį dvinario daugiklį reiškinyje $ab(abc-1)+bc(abc-1)-ca(abc-1)$.
        \item Raskite bendrąjį dvinario daugiklį reiškinyje $x(x-x_1)-x_2(x-x_1)$.
        \item Raskite bendrąjį dvinario daugiklį, atsižvelgdami į $x_1-x=-(x-x_1)$.
    \end{enumerate}
    \item Sugrupuokite narius į dvi poras ir kiekvienoje poroje iškelkite bendrąjį daugiklį.
    \begin{enumerate}[label={(\arabic*)}]
        \item Sugrupuokite narius poromis: $(x^2+2x)+(3x+6)$.
        \item Sugrupuokite narius poromis: $(2a^2-6a)+(5a-15)$.
        \item Sugrupuokite narius poromis: $(6m^2+2m)+(9m+3)$.
        \item Sugrupuokite narius poromis: $(y^2-4y)+(-y+4)$.
        \item Sugrupuokite narius poromis: $(x^2-7x)+(-2x+14)$.
        \item Sugrupuokite narius poromis: $(ab^2+abc)+(bc+c^2)$.
        \item Sugrupuokite narius poromis: $(x^3+3x^2)+(2x+6)$.
        \item Sugrupuokite narius poromis: $(x^3-x^2)+(-y^2x+y^2)$.
    \end{enumerate}
    \item Iškelkite $a$ prieš skliaustus ir prisiminkite Vijeto teoremą.
    \item Raskite šaknis pagal Vijeto teoremą arba diskriminanto formulę.

    \item Išskaidykite skaitiklį ir vardiklį dauginamaisiais; pasižymėkite, kurioms $x$ reikšmėms vardiklis lygus nuliui.
    \begin{enumerate}[label={(\arabic*)}]
        \item Išskaidykite kvadratinį trinarį skaitiklyje.
        \item Skaitiklyje panaudokite kvadratų skirtumo formulę; išskaidykite vardiklio trinarį.
        \item Išskaidykite skaitiklio trinarį, atsižvelgdami į tai, kad jo pirmojo nario koeficientas yra $2$.
        \item Iškelkite $3x$ skaitiklyje ir išskaidykite vardiklį.
        \item Išskaidykite kvadratinį trinarį skaitiklyje.
        \item Skaitiklyje panaudokite kvadratų skirtumo formulę ir išskaidykite vardiklį.
    \end{enumerate}
\end{enumerate}

\newpage
\section{Hintai 2}

\begin{enumerate}[label=\textbf{\arabic*.}, leftmargin=*]
    \item Iškelkite bendrąjį daugiklį ir užrašykite likusius narius skliaustuose.
    \begin{enumerate}[label={(\arabic*)}]
        \item $3x+12=3\cdot x+3\cdot4$.
        \item $5ab+10ac=5a\cdot b+5a\cdot2c$.
        \item $x^2+7x=x\cdot x+x\cdot7$.
        \item $4x^2-12x=4x\cdot x+4x\cdot(-3)$.
        \item $-2x+10=(-2)\cdot x+(-2)\cdot(-5)$.
        \item $-x_2x+x_1x_2=x_2\cdot(-x)+x_2\cdot x_1$.
    \end{enumerate}
    \item Iškelkite bendrąjį dvinario daugiklį; prireikus pakeiskite reiškinio ženklą.
    \begin{enumerate}[label={(\arabic*)}]
        \item $3x(x+2)+5(x+2)=3x(x+2)+5\cdot1(x+2)$.
        \item $y(2y-1)-4(2y-1)=y(2y-1)+(-4)(2y-1)$.
        \item $a(a-3)+(a-3)=a(a-3)+1(a-3)$.
        \item $2m(m+4)-3n(m+4)=2m(m+4)+(-3n)(m+4)$.
        \item $x(x-5)-(x-5)=x(x-5)+(-1)(x-5)$.
        \item $-3x(x-4)-2(4-x)=-3x(x-4)+2(x-4)$.
        \item $ab(b-2c)-c(2c-b)=ab(b-2c)+c(b-2c)$.
        \item $a(b-a)+b(a-b)-c(a-b)=a(b-a)-b(b-a)+c(b-a)$.
        \item $x(y-z)+y(z-y)-xyz(y-z)=x(y-z)-y(y-z)-xyz(y-z)$.
        \item $ab(abc-1)+bc(abc-1)-ca(abc-1)$.
        \item $x(x-x_1)-x_2(x-x_1)=x(x-x_1)+(-x_2)(x-x_1)$.
        \item $x(x-x_1)-x_2(x_1-x)=x(x-x_1)+x_2(x-x_1)$.
    \end{enumerate}
    \item Išskaidę abi poras, patikrinkite, ar jose atsirado tas pats dvinario daugiklis.
    \begin{enumerate}[label={(\arabic*)}]
        \item $x(x+2)+3(x+2)$.
        \item $2a(a-3)+5(a-3)$.
        \item $2m(3m+1)+3(3m+1)$.
        \item $y(y-4)-1(y-4)$.
        \item $x(x-7)-2(x-7)$.
        \item $ab(b+c)+c(b+c)$.
        \item $x^2(x+3)+2(x+3)$.
        \item $x^2(x-1)-y^2(x-1)$.
    \end{enumerate}
    \item Pagal Vijeto teoremą: $\frac{b}{a}=-(x_1+x_2)$, $\frac{c}{a}=x_1\cdot x_2$.
    \item Pritaikykite šaknų formulę $a(x-x_1)(x-x_2)$.
    \begin{enumerate}[label={(\arabic*)}]
        \item Šaknys yra $2$ ir $4$.
        \item Šaknys yra $\frac16$ ir $-1$.
        \item Šaknys yra $-3$ ir $4$.
        \item Šaknys yra $\frac15$ ir $-\frac12$.
        \item Šaknys yra $2$ ir $7$.
        \item Šaknys yra $3$ ir $-6$.
        \item Šaknys yra $3$ ir $4$.
        \item Šaknys yra $1$ ir $-3$.
        \item Šaknys yra $\frac12$ ir $-3$.
        \item Šaknys yra $1$ ir $-\frac25$.
        \item Šaknys yra $3$ ir $-\frac23$.
        \item Šaknys yra $\frac32$ ir $-1$.
    \end{enumerate}
    \item Išskaidę skaitiklį ir vardiklį,  suprastinkite bendrąjį daugiklį.
    \begin{enumerate}[label={(\arabic*)}]
        \item Skaitiklis yra $(x-5)(x+2)$, vardiklis — $(x-5)$.
        \item Skaitiklis yra $(x-3)(x+3)$, vardiklis — $(x-3)(x+2)$.
        \item Skaitiklis yra $(2x-1)(x+1)$, vardiklis — $(x+1)$.
        \item Skaitiklis yra $3x(x-4)$, vardiklis — $(x-4)(x+2)$.
        \item Skaitiklis yra $(x-6)(x+2)$, vardiklis — $(x-6)$.
        \item Skaitiklis yra $(x-2)(x+2)$, vardiklis — $(x-2)(x+3)$.
    \end{enumerate}
\end{enumerate}

\newpage
\section{Hintai 3}

\begin{enumerate}[label=\textbf{\arabic*.}, leftmargin=*]
    \item Patikrinkite iškeliamąjį daugiklį išskliausdami gautą sandaugą.
    \begin{enumerate}[label={(\arabic*)}]
        \item $3x+12=3(x+4)$.
        \item $5ab+10ac=5a(b+2c)$.
        \item $x^2+7x=x(x+7)$.
        \item $4x^2-12x=4x(x-3)$.
        \item $-2x+10=-2(x-5)$.
        \item $-x_2x+x_1x_2=x_2(x_1-x)$.
    \end{enumerate}
    \item Iškelkite bendrą dvinario daugiklį ir sutvarkykite likusią skliaustų dalį.
    \begin{enumerate}[label={(\arabic*)}]
        \item $3x(x+2)+5(x+2)=(x+2)(3x+5)$.
        \item $y(2y-1)-4(2y-1)=(2y-1)(y-4)$.
        \item $a(a-3)+(a-3)=(a-3)(a+1)$.
        \item $2m(m+4)-3n(m+4)=(m+4)(2m-3n)$.
        \item $x(x-5)-(x-5)=(x-5)(x-1)$.
        \item $-3x(x-4)-2(4-x)=(x-4)(2-3x)$.
        \item $ab(b-2c)-c(2c-b)=(b-2c)(ab+c)$.
        \item $a(b-a)+b(a-b)-c(a-b)=(a-b)(b-a-c)$.
        \item $x(y-z)+y(z-y)-xyz(y-z)=(y-z)(x-y-xyz)$.
        \item $ab(abc-1)+bc(abc-1)-ca(abc-1)=(abc-1)(ab+bc-ca)$.
        \item $x(x-x_1)-x_2(x-x_1)=(x-x_1)(x-x_2)$.
        \item $x(x-x_1)-x_2(x_1-x)=(x-x_1)(x+x_2)$.
    \end{enumerate}
    \item Iškelkite porose bendruosius dvinario daugiklius, tada iškelkite juos iš visos išraiškos.
    \begin{enumerate}[label={(\arabic*)}]
        \item $(x+2)(x+3)$.
        \item $(a-3)(2a+5)$.
        \item $(3m+1)(2m+3)$.
        \item $(y-4)(y-1)$.
        \item $(x-7)(x-2)$.
        \item $(b+c)(ab+c)$.
        \item $(x+3)(x^2+2)$.
        \item Iškelkite bendrąjį daugiklį, tada pritaikykite kvadratų skirtumo formulę: $$(x-1)(x^2-y^2)=(x-1)(x-y)(x+y)$$.
    \end{enumerate}
    \item $ax^2+bx+c=a(x^2-x_1x-x_2x+x_1x_2)$.
    \item Išskaidę trinarį gaunate:
    \begin{enumerate}[label={(\arabic*)}]
        \item $x^2-6x+8=(x-2)(x-4)$.
        \item $6x^2+5x-1=(6x-1)(x+1)$.
        \item $x^2-x-12=(x+3)(x-4)$.
        \item $10x^2+3x-1=(5x-1)(2x+1)$.
        \item $x^2-9x+14=(x-7)(x-2)$.
        \item $-x^2-3x+18=(3-x)(x+6)$.
        \item $2x^2-14x+24=2(x-3)(x-4)$.
        \item $-3x^2-6x+9=-3(x-1)(x+3)$.
        \item $2x^2+5x-3=(2x-1)(x+3)$.
        \item $-5x^2+3x+2=(1-x)(5x+2)$.
        \item $3x^2-7x-6=(3x+2)(x-3)$.
        \item $-2x^2+x+3=(3-2x)(x+1)$.
    \end{enumerate}
    \item Suprastinkite bendruosius daugiklius gausite:
    \begin{enumerate}[label={(\arabic*)}]
        \item $\frac{(x-5)(x+2)}{x-5}=x+2$, kai $x\ne5$.
        \item $\frac{(x-3)(x+3)}{(x-3)(x+2)}=\frac{x+3}{x+2}$, kai $x\ne3,-2$.
        \item $\frac{(2x-1)(x+1)}{x+1}=2x-1$, kai $x\ne-1$.
        \item $\frac{3x(x-4)}{(x-4)(x+2)}=\frac{3x}{x+2}$, kai $x\ne4,-2$.
        \item $\frac{(x-6)(x+2)}{x-6}=x+2$, kai $x\ne6$.
        \item $\frac{(x-2)(x+2)}{(x-2)(x+3)}=\frac{x+2}{x+3}$, kai $x\ne2,-3$.
    \end{enumerate}
\end{enumerate}

\end{document}
